\chapter{黎曼度量与Levi--Civita联络}
\label{ch:11-riemannian}
\ChapterPlan{周4.1及4.3（本地页码待核对）；梁上册第二版3.2.2、3.4及第3章习题7、12、13、17；度量存在与截面曲率另作补充}{正文及证明初稿；附8题详解与3幅原创图，尚未独立数学审阅}

第10章允许任意联络，因而平行移动不必保持长度。本章先在每个切空间选内积，再要求求导与内积相容；加上无挠条件，便从许多联络中唯一选出Levi--Civita联络。随后把联络曲率降为四阶协变张量，用其代数对称性理解不同的曲率缩并。

全章流形均光滑、Hausdorff且第二可数，无边界，度量正定；积分公式另说明定向与紧支条件。保持第10章的方向槽位及$R(X,Y)Z=\nabla_X\nabla_YZ-\nabla_Y\nabla_XZ-\nabla_{[X,Y]}Z$。为避免降指标后的数组含义不明，专记
\[
 \mathcal R(X,Y,Z,W)=g(R(X,Y)Z,W).
\]
在此约定下，截面曲率的分子为$\mathcal R(u,v,v,u)$，球面取正曲率。这与第2章的$K=\det S$一致，不用改变形算子的符号。

\Needspace{17\baselineskip}
\section{黎曼度量与度量存在性}
\begin{definition}[黎曼度量]\label{def:riemannian-metric}
$M$上的黎曼度量是光滑对称$(0,2)$型张量场$g$，使每个$g_p$正定，即$v\ne0$时$g_p(v,v)>0$。称$(M,g)$为黎曼流形，记$|v|_g=\sqrt{g(v,v)}$。非零向量夹角由$\cos\theta=g(u,v)/(|u|_g|v|_g)$定义；Cauchy--Schwarz不等式保证其值在$[-1,1]$中。
\end{definition}

在坐标卡中，$g=g_{ij}\dd x^i\otimes\dd x^j$，$(g_{ij})$为光滑正定对称矩阵。换坐标$y=y(x)$时
\[
 g'_{ab}=g_{ij}\frac{\partial x^i}{\partial y^a}\frac{\partial x^j}{\partial y^b},
 \qquad (g'_{ab})=J^{\mathsf T}(g_{ij})J,\quad J=\frac{\partial x}{\partial y}.
\]
这里两个下槽均按协变规律改变。矩阵元素不能在每张任意坐标卡中都指定为$\delta_{ij}$；必须满足交叠相容式。每点可选正交归一基，只是逐纤维的线性代数结论。

\begin{theorem}[度量的存在与凸性]\label{thm:metric-existence}
每个本章意义下的光滑流形都有黎曼度量。若$g,h$为两个度量，$0\leq t\leq1$，则$g_t=(1-t)g+th$是光滑一族度量；更一般，度量的局部有限非负加权和在每点至少有一个正权重时仍正定。
\end{theorem}
\begin{proof}
取坐标覆盖$\{U_\alpha\}$及第8章构造的光滑单位分解$\{\rho_\alpha\}$，满足$\operatorname{supp}\rho_\alpha\subset U_\alpha$。在$U_\alpha$拉回欧氏内积，得到$g_\alpha=\sum_i\dd x_\alpha^i\otimes\dd x_\alpha^i$。令
\[
 g=\sum_\alpha\rho_\alpha g_\alpha,
\]
其中每项在卡外延为零。支集条件保证任意卡外点的某个邻域上该项恒零，所以零延拓光滑；局部有限性使总和光滑。固定$p$及$0\ne v\in T_pM$，所有非零权重项的$g_\alpha(v,v)$严格正，而权重和为1，故$g_p(v,v)>0$。同样的逐点论证给出加权和结论及$g_t$的正定性，参数光滑性直接可见。
\end{proof}

这一构造不选定全局标架，也不要求流形可定向。正定矩阵的凸性是关键；固定不定符号的非退化矩阵集合一般不凸，因此不能照此声称每个流形都有指定符号的伪黎曼度量。

\Needspace{17\baselineskip}
\begin{proposition}[局部正交标架与升降指标]\label{prop:orthonormal-musical}
每点邻域上都有光滑$g$正交归一标架。映射$\flat:TM\to T^*M$、$v^\flat=g(v,\cdot)$为光滑丛同构，其逆记$\sharp$。在坐标中
\[
 (v^\flat)_i=g_{ij}v^j,\qquad (\alpha^\sharp)^i=g^{ij}\alpha_j,
 \qquad(g^{ij})=(g_{ij})^{-1}.
\]
对$f\in C^\infty(M)$，梯度$\operatorname{grad}_g f=(\dd f)^\sharp$由$g(\operatorname{grad}_g f,X)=Xf$唯一确定。
\end{proposition}
\begin{proof}
从坐标基逐步作Gram--Schmidt：$u_k=\partial_k-\sum_{a<k}g(\partial_k,e_a)e_a$，$e_k=u_k/\sqrt{g(u_k,u_k)}$。线性无关性保证每步分母正，正函数的平方根光滑，所以所得标架光滑。逆矩阵由伴随矩阵除以正行列式给出，其系数光滑；非退化性保证$\flat$可逆。梯度结论是把协向量$\dd f$代入此同构。
\end{proof}

正交标架一般不是坐标基。例如平面极坐标的$e_1=\partial_r$、$e_2=r^{-1}\partial_\theta$满足$[e_1,e_2]=-r^{-1}e_2$，所以即使度量在此标架的矩阵恒为$I$，也不能据此断言联络矩阵为零。

\begin{figure}[htbp]
\centering\DGDiagram{metric-unit-balls}

\caption{同一二维向量空间的两种内积。采用相同坐标比例，$g_0=I$的单位圆与$g_1=\operatorname{diag}(4,1)$的单位椭圆不同；这不是流形在环境空间中的形状。}
\label{fig:metric-unit-balls}
\end{figure}

\Needspace{17\baselineskip}
\begin{proposition}[拉回度量]\label{prop:pullback-metric}
若$F:M\to N$光滑且$h$为$N$的黎曼度量，则
\[
 (F^*h)_p(u,v)=h_{F(p)}(\dd F_pu,\dd F_pv)
\]
为光滑半正定对称张量；它正定当且仅当$F$为浸入。若$F$为微分同胚且$F^*h=g$，称$F$为等距映射；若仅为局部微分同胚并满足此式，称局部等距映射。
\end{proposition}
\begin{proof}
光滑性与对称性来自张量拉回。$h(\dd Fv,\dd Fv)=0$恰当$\dd Fv=0$，所以逐点正定等价于$\ker\dd F_p=0$。这也解释为何任意光滑映射的拉回未必是度量：常映射的拉回恒为零。
\end{proof}

在$\R^3$中半径$a>0$的球面取包含映射的拉回度量，球坐标卡$0<\theta<\pi$上为$a^2(\dd\theta^2+\sin^2\theta\,\dd\phi^2)$。柱面取$a^2\dd\theta^2+\dd z^2$，局部坐标$u=a\theta$使它成为$\dd u^2+\dd z^2$；球面与柱面的内蕴曲率不能由两者都在三维空间中弯曲来判断。长度、长度距离及其拓扑和极值性质在第12章系统建立。


\Needspace{17\baselineskip}
\section{Koszul公式与基本定理}
\begin{definition}[度量相容]
切丛联络$\nabla$与$g$相容，是指第10章诱导张量联络满足$\nabla g=0$，即
\[
 Xg(Y,Z)=g(\nabla_XY,Z)+g(Y,\nabla_XZ).
\]
坐标式为$\partial_i g_{jk}-\Gamma^a_{ij}g_{ak}-\Gamma^a_{ik}g_{ja}=0$。相容性与无挠是两个不同条件。
\end{definition}

\begin{theorem}[黎曼几何基本定理与Koszul公式]\label{thm:levi-civita}
给定黎曼度量$g$，存在唯一无挠且与$g$相容的切丛联络，称Levi--Civita联络。它由以下Koszul公式确定：
\begin{align*}
 2g(\nabla_XY,Z)={}&Xg(Y,Z)+Yg(Z,X)-Zg(X,Y)\\
 &-g(X,[Y,Z])+g(Y,[Z,X])+g(Z,[X,Y]).
\end{align*}
\end{theorem}
\begin{proof}
先证必要性。写出相容性对$(X,Y,Z)$、$(Y,Z,X)$及$(Z,X,Y)$的三个式子，将前两式相加并减第三式。用$\nabla_XY-\nabla_YX=[X,Y]$消去反向导数，恰得上式。因此任何合格联络的$\nabla_XY$与所有$Z$的内积已确定；$g$非退化，故唯一。

再证存在。记右端为$K(X,Y,Z)$，不预先假定它来自联络。它对各槽实线性。对光滑函数$f$直接代入括号的Leibniz式，得到
\begin{align*}
 K(X,Y,fZ)&=fK(X,Y,Z),\\
 K(fX,Y,Z)&=fK(X,Y,Z),\\
 K(X,fY,Z)&=fK(X,Y,Z)+2X(f)g(Y,Z).
\end{align*}
第一式中额外项依次为$X(f)g(Y,Z)+Y(f)g(X,Z)-Y(f)g(X,Z)-X(f)g(Y,Z)=0$。第二式中$Y(f)g(Z,X)-Z(f)g(X,Y)$与两个括号的相反项抵消；第三式中$Y$槽带来的导数相消后只留下两份$X(f)g(Y,Z)$。因此$K(X,Y,\cdot)/2$是光滑一形式，经$\sharp$定义光滑向量场$\nabla_XY$。后两式给出联络的函数线性与Leibniz规则，故确为全局联络。

交换前两槽，直接整理得
\[
 K(X,Y,Z)-K(Y,X,Z)=2g([X,Y],Z),
\]
所以$T=0$。交换后两槽并相加，得到
\[
 K(X,Y,Z)+K(X,Z,Y)=2Xg(Y,Z),
\]
所以$\nabla g=0$。这完成存在性，不需要预先挑选某个全局坐标系来拼接系数。
\end{proof}

\begin{corollary}[Christoffel公式]\label{cor:metric-christoffel}
在坐标基中，Levi--Civita联络满足
\[
 \Gamma^k_{ij}=\frac12g^{k\ell}
 (\partial_i g_{j\ell}+\partial_j g_{i\ell}-\partial_\ell g_{ij}).
\]
\end{corollary}
\begin{proof}
Koszul公式中的坐标括号都为零，故$2g_{k\ell}\Gamma^k_{ij}=\partial_i g_{j\ell}+\partial_j g_{i\ell}-\partial_\ell g_{ij}$，乘逆矩阵即可。公式只在坐标基中如此简化；一般正交标架须保留括号项。
\end{proof}

例如$g=\dd r^2+r^2\dd\theta^2$给出$\Gamma^r_{\theta\theta}=-r$、$\Gamma^\theta_{r\theta}=\Gamma^\theta_{\theta r}=1/r$，与第10章通过坐标变换得到的系数相同。非零系数不代表曲率非零。

\Needspace{17\baselineskip}
\begin{proposition}[相容性与等距平行移动]\label{prop:metric-transport}
对任意切丛联络，$\nabla g=0$当且仅当沿每条光滑曲线的平行移动保持内积。相容联络沿曲线满足
\[
 \frac{\dd}{\dd t}g(V,W)=g(D_tV,W)+g(V,D_tW).
\]
特别地，Levi--Civita仿射测地线的速度范数为常数。
\end{proposition}
\begin{proof}
对沿曲线任意$V,W$在坐标中求导，减去右端，剩余正是$(\nabla_{\dot\gamma}g)(V,W)$，所以相容时式子成立。平行场使右端为零，内积保持。

反之，在$p$取任意$u,v,w\in T_pM$，选曲线满足$\dot\gamma(0)=u$，沿其平行延拓$v,w$为$V,W$。由假设$g(V,W)$恒定，前述坐标计算得$(\nabla_u g)(v,w)=0$。三个向量任意，所以$\nabla g=0$。最后对$V=W=\dot\gamma$用$D_t\dot\gamma=0$即得速度平方导数为零；常值曲线也包含在内。
\end{proof}

\begin{figure}[htbp]
\centering\DGDiagram{metric-transport}

\caption{在起终点分别用正交归一基表示二维切空间。相容联络的平行移动同时保持两向量的长度与夹角，图中两单位圆表示不同纤维；图示不声称有沿所有路径一致的全局识别。}
\label{fig:metric-transport}
\end{figure}

\begin{proposition}[等距映射、常数缩放与乘积]\label{prop:metric-naturality}
局部等距映射保持Levi--Civita联络及其曲率，因而保持带仿射参数的测地线。正数常量$c$的缩放$\widetilde g=cg$不改变Levi--Civita联络及$(1,3)$型曲率。乘积度量$g_1\oplus g_2$的联络在来自两个因子的提升场上分别为各因子的联络，混合导数为零。
\end{proposition}
\begin{proof}
先将局部等距映射限制为某邻域的微分同胚，把目标联络拉回：
\[
 \widehat\nabla_XY=(\dd F)^{-1}(\nabla^N_{F_*X}F_*Y).
\]
它满足联络规则；相容性由$F^*h=g$，无挠性由$F_*[X,Y]=[F_*X,F_*Y]$给出。因此由唯一性$\widehat\nabla=\nabla^M$。代入曲率和沿曲线导数的定义得其余保持性质。

常数缩放时Koszul等式两边同乘$c$，故联络相同；曲率的定义只涉及联络，因而相同。对乘积，提升场来自不同因子时括号为零，混合内积为零，同因子内积对另一因子求导也为零。代入Koszul得到混合导数为零、同因子导数为原联络。一般场再用联络规则从提升局部标架扩展，不能误把依赖两组变量的一般场的系数导数删掉。
\end{proof}

\subsection*{体积与标量微分算子}
若$M$已定向，令$d=\det(g_{ij})>0$。在正向卡中定义$\mathrm{vol}_g=\sqrt d\,\dd x^1\wedge\cdots\wedge\dd x^n$。换卡时$d'=(\det J)^2d$，而正向换卡的Jacobian为正，故这些形式拼成全局正体积形式；不可定向时仍有局部密度，但本章不将其当全局顶次形式积分。

\begin{proposition}[散度与Laplace算子]\label{prop:metric-divergence}
记$\operatorname{div}_gX=\operatorname{tr}(Y\mapsto\nabla_YX)$，取符号约定$\Delta_g f=\operatorname{div}_g\operatorname{grad}_g f$。则
\begin{align*}
 \Gamma^i_{ik}&=\partial_k\log\sqrt d,\\
 \operatorname{div}_gX&=\frac1{\sqrt d}\partial_i(\sqrt d\,X^i),\\
 \Delta_g f&=\frac1{\sqrt d}\partial_i(\sqrt d\,g^{ij}\partial_j f)
 =g^{ij}(\partial_i\partial_j f-\Gamma^k_{ij}\partial_k f).
\end{align*}
在定向流形上还满足$\dd(\iota_X\mathrm{vol}_g)=(\operatorname{div}_gX)\mathrm{vol}_g$。若$f,h$光滑且至少一个紧支，则
\[
 \int_M f\Delta_g h\,\mathrm{vol}_g
 =-\int_M g(\operatorname{grad}_g f,\operatorname{grad}_g h)\,\mathrm{vol}_g.
\]
\end{proposition}
\begin{proof}
在Christoffel公式中缩并上下$i$，第一、第三项重命名哑指标后抵消，得$\Gamma^i_{ik}=\tfrac12g^{ij}\partial_k g_{ij}$。行列式按列求导，余子式公式给$\partial_k d=d\,g^{ij}\partial_k g_{ji}$，于是得到第一式。取$\nabla X$的迹并用乘积法则即得第二式；代入梯度分量得第三式前半。又因$\nabla g=0$，$\nabla\operatorname{grad}h$的迹等于$\nabla\dd h$的度量迹，从第10章Hessian分量式得到后半。

对$\iota_X\mathrm{vol}_g$逐项外微分，只有缺少的$\dd x^i$方向存活，各项符号抵消后给
\[
 \partial_i(\sqrt d X^i)\dd x^1\wedge\cdots\wedge\dd x^n.
\]
再取$X=f\operatorname{grad}h$，用散度的乘积规则得
\[
 \operatorname{div}(f\operatorname{grad}h)
 =g(\operatorname{grad}f,\operatorname{grad}h)+f\Delta h.
\]
$f$紧支时$X$紧支，$h$紧支时$\operatorname{grad}h$紧支，故两种情形均可用第8章紧支Stokes，得到积分等式。
\end{proof}

本书的$\Delta$在欧氏卡中是$\sum_i\partial_i^2$，所以$\int f\Delta f\leq0$。谱论中也常把其负号称为Laplace算子，比较公式时须先核对约定。这里尚未证明任何椭圆估计或Hodge分解。


\Needspace{17\baselineskip}
\section{曲率张量的对称性}
\begin{theorem}[Levi--Civita曲率的代数对称性]\label{thm:metric-curvature-symmetries}
令$\mathcal R(X,Y,Z,W)=g(R(X,Y)Z,W)$。则
\begin{align*}
 \mathcal R(X,Y,Z,W)&=-\mathcal R(Y,X,Z,W),\\
 \mathcal R(X,Y,Z,W)&=-\mathcal R(X,Y,W,Z),\\
 \mathcal R(X,Y,Z,W)&=\mathcal R(Z,W,X,Y),\\
 \mathcal R(X,Y,Z,W)+\mathcal R(Y,Z,X,W)+\mathcal R(Z,X,Y,W)&=0.
\end{align*}
后两式分别称成对交换对称性与代数Bianchi恒等式。
\end{theorem}
\begin{proof}
第一式由曲率定义。由度量相容，诱导张量联络的曲率作用于$g$为零；第10章对偶与张量曲率公式给
\[
 0=(R^{T^*\otimes T^*}(X,Y)g)(Z,W)
 =-g(R(X,Y)Z,W)-g(Z,R(X,Y)W),
\]
这就是第二式。由无挠及第10章第一Bianchi，最后一式成立。

成对交换只需代数，不用法坐标。固定四向量，记$A_{abcd}=\mathcal R(a,b,c,d)$。先用Bianchi和末二槽反对称性，再分别对$(b,c,d)$及$(c,a,d)$用Bianchi，得
\begin{align*}
 A_{abcd}&=A_{bcda}+A_{cadb}\\
 &=2A_{cdab}+A_{dbac}+A_{adbc}.
\end{align*}
另一方面对$(d,b,a)$用Bianchi，$A_{dbac}+A_{badc}+A_{adbc}=0$，而两次反对称给$A_{badc}=A_{abcd}$。因此$A_{abcd}=2A_{cdab}-A_{abcd}$，即成对交换对称性。
\end{proof}

这些假设各有作用：对任意联络只有方向反对称总成立；$\nabla g=0$给末二槽反对称；无挠给右端为零的代数Bianchi；成对交换用到了后两者。相容但有挠的联络不能直接套用全套对称性。以正交归一标架记$\mathcal R_{ijkl}=g(R(e_i,e_j)e_k,e_l)$时，球面的$\mathcal R_{1221}$为正，$\mathcal R_{1212}$则为负。

\begin{proposition}[截面二次值确定整个曲率]\label{prop:curvature-polarization}
若四线性张量$A$满足上述代数对称性，且对任意$x,y$有$A(x,y,y,x)=0$，则$A=0$。因而同一内积空间上两个曲率张量若对所有二维平面有相同截面曲率，就完全相同。
\end{proposition}
\begin{proof}
先将$x$替成$x+z$，减去各自的对角式。由成对交换及首末反对称，两个交叉项相等，故$A(x,y,y,z)=0$。再将$y$替成$y+w$得到
\[
 A(x,y,w,z)+A(x,w,y,z)=0.
\]
所以$A$对第二、第三槽也反对称；连同首二槽反对称，前三槽全交替。循环置换三个槽为偶置换，代数Bianchi于是成为$3A(x,y,w,z)=0$，得$A=0$。比较两个张量时取差；线性相关的$x,y$由反对称自动给零，独立向量的分子由截面曲率相等给零，故适用前述论证。
\end{proof}

\begin{proposition}[黎曼平坦的局部含义]\label{prop:riemannian-flat}
黎曼度量的曲率为零，当且仅当每点附近有坐标使$g_{ij}=\delta_{ij}$。等价地，它局部等距于欧氏空间的开集。
\end{proposition}
\begin{proof}
若$R=0$，第10章定理\ref{thm:flat-parallel-frame}给局部平行标架。先在一点选正交归一初始基，再作该构造；因$\nabla g=0$，平行标架的内积沿任意方向导数均为零，缩小到连通邻域后恒为$\delta_{ab}$。Levi--Civita无挠，所以$[e_a,e_b]=\nabla_{e_a}e_b-\nabla_{e_b}e_a=0$；第5章交换流构造坐标，使$e_a=\partial/\partial y^a$。于是在此坐标中$g_{ab}=\delta_{ab}$。反之度量坐标分量常为$I$时，Christoffel及曲率分量均为零。
\end{proof}

这一结论是邻域上的平坦性，不是只在一点让系数变简单。法坐标可在一点令$g_{ij}=\delta_{ij}$、$\Gamma^k_{ij}=0$，但二阶度量变化仍可携带曲率；法坐标的指数映射构造留第12章。


\Needspace{17\baselineskip}
\section{截面曲率、Ricci曲率与标量曲率}
\begin{definition}[截面曲率]\label{def:sectional-curvature}
设$n\geq2$，$\sigma\subset T_pM$为二维线性子空间，取任意基$u,v$，定义
\[
 K_p(\sigma)=\frac{g(R(u,v)v,u)}{g(u,u)g(v,v)-g(u,v)^2}.
\]
分母是正的Gram行列式，即该平行四边形的面积平方。$K$与$\sigma$的定向无关。
\end{definition}
\begin{proof}[良定义性]
若$u'=au+bv$、$v'=cu+dv$，$ad-bc\ne0$，方向反对称给$R(u',v')=(ad-bc)R(u,v)$，末二槽反对称又给
\[
 g(R(u',v')v',u')=(ad-bc)^2g(R(u,v)v,u).
\]
Gram行列式在同一换基下也乘$(ad-bc)^2$，故比值不变。特别地，正交归一基下$K(\sigma)=\mathcal R(e_1,e_2,e_2,e_1)$。
\end{proof}

\begin{proposition}[常截面值的代数公式]\label{prop:constant-curvature-model}
固定$p$，所有二维平面的截面曲率都等于$k$，当且仅当
\[
 R(u,v)w=k\bigl(g(v,w)u-g(u,w)v\bigr).
\]
这里$k$先是该点的数；各点相同截面值时可写为函数$k(p)$，其全局常数性尚须证明。
\end{proposition}
\begin{proof}
右侧降指标后满足两种反对称、成对交换及代数Bianchi（逐项两两抵消），代入$(u,v,v,u)$得$k$乘Gram行列式，因而截面值为$k$。反之，用真实曲率减去此模型张量，所有截面二次值为零，命题\ref{prop:curvature-polarization}说明其差为零。
\end{proof}

\Needspace{17\baselineskip}
\begin{definition}[Ricci曲率与标量曲率]\label{def:ricci-scalar}
定义$\Ric(X,Y)=\operatorname{tr}(Z\mapsto R(Z,X)Y)$，再用度量取迹得$\mathrm{Scal}=\operatorname{tr}_g\Ric$。坐标及正交归一标架表达为
\begin{align*}
 \Ric_{jk}&=R^i{}_{kij},&\quad \mathrm{Scal}&=g^{jk}\Ric_{jk},\\
 \Ric(X,Y)&=\sum_a g(R(e_a,X)Y,e_a),&
 \mathrm{Scal}&=\sum_{a,b}\mathcal R(e_a,e_b,e_b,e_a).
\end{align*}
Ricci定义中的第一个迹是线性变换的迹，不先要求度量；但此处的对称性及后续截面解释使用Levi--Civita条件。
\end{definition}

\begin{proposition}[缩并、方向平均与缩放]\label{prop:ricci-traces}
$\Ric$为光滑对称$(0,2)$型张量。若$|u|=1$，补成正交归一基$e_1=u,e_2,\ldots,e_n$，则
\[
 \Ric(u,u)=\sum_{a=2}^n K(\operatorname{span}\{u,e_a\}),\qquad
 \mathrm{Scal}=2\sum_{a<b}K(\operatorname{span}\{e_a,e_b\}).
\]
若该点截面曲率均为$k$，则$\Ric=(n-1)kg$、$\mathrm{Scal}=n(n-1)k$。对常数缩放$\widetilde g=cg$，
\[
 \widetilde K=c^{-1}K,\quad \widetilde\Ric=\Ric,\quad
 \widetilde{\mathrm{Scal}}=c^{-1}\mathrm{Scal},\quad \widetilde{\mathcal R}=c\mathcal R.
\]
\end{proposition}
\begin{proof}
曲率的张量性与迹的基无关性给$\Ric$的型别及光滑性。代数Bianchi给
\[
 \mathcal R(e_a,X,Y,e_a)-\mathcal R(e_a,Y,X,e_a)
 =-\mathcal R(X,Y,e_a,e_a)=0,
\]
逐项相等，故对称。代入正交归一基，$a=1$项为零，其余各项恰为对应截面曲率；再对$u=e_b$求和，每个无序对出现两次。模型公式直接取迹给常截面值情形。

缩放时$(1,3)$型$R$不变，故降指标的$\mathcal R$乘$c$，Gram行列式乘$c^2$，$K$除以$c$。Ricci为不依赖度量的同一变换之迹，不变；标量再乘逆度量，故除以$c$。特别注意“Ricci不变”说的是$(0,2)$张量，不是其相对新度量的特征值不变。
\end{proof}

\begin{figure}[htbp]
\centering\DGDiagram{sectional-trace}

\caption{三维切空间中固定单位方向$e_1$，Ricci值由两个包含它的正交截面曲率相加得到。图为线性空间的投影，不代表流形上两张实际曲面。一般$n$维为$n-1$项之和。}
\label{fig:sectional-trace}
\end{figure}

\subsection*{二维、模型与乘积}
当$n=2$时，每点只有一个二维平面，记其曲率为$K(p)$。模型公式逐点成立，从而
\[
 \Ric=Kg,\qquad \mathrm{Scal}=2K,\qquad
 G:=\Ric-\tfrac12\mathrm{Scal}\,g=0.
\]
这不表示$K=0$，也不表示$K$恒定。对欧氏空间中的曲面，第2章的切向投影联络满足内积乘法式（法向部分与切向量正交），且由环境求导的$D_XY-D_YX=[X,Y]$可知无挠；基本定理说明它正是诱导度量的Levi--Civita联络。第2章Gauss方程给$g(R(u,v)v,u)=II(u,u)II(v,v)-II(u,v)^2$，故此处的$K$正好是$\det S$。半径$a$的球面因此$K=a^{-2}$、$\Ric=a^{-2}g$、$\mathrm{Scal}=2a^{-2}$；习题将独立用坐标验证，不借用第13章的一般子流形公式。平面及圆柱面$K=0$，圆柱面的外在弯曲不改变此结论。

乘积度量的曲率在因子方向为因子曲率，混合方向为零，这是把命题\ref{prop:metric-naturality}的联络代入曲率定义的直接结果。例如$S^2(a)\times\R$中，球面切平面的$K=a^{-2}$，球面与实线各取一向量的平面$K=0$，因此该空间没有常截面曲率。相对正交基，Ricci特征值为$(a^{-2},a^{-2},0)$，标量为$2a^{-2}$；一个标量不足以保存各个平面方向的信息。当$n=1$时曲率因前两槽反对称恒零，Ricci和标量均零，但没有二维截面可供定义$K$。

\Needspace{17\baselineskip}
\begin{theorem}[缩并Bianchi与Einstein张量]\label{thm:contracted-bianchi}
Levi--Civita联络满足
\[
 \nabla^i\Ric_{ij}=\frac12\partial_j\mathrm{Scal},\qquad
 \nabla^iG_{ij}=0,\quad G=\Ric-\frac12\mathrm{Scal}\,g.
\]
这里$\nabla^i=g^{ik}\nabla_k$，协变导数包含所有张量槽的联络项。
\end{theorem}
\begin{proof}
写$A_{abcd}=\mathcal R_{abcd}$。无挠的第二Bianchi与$\nabla g=0$给
\[
 \nabla_a A_{bcde}+\nabla_b A_{cade}+\nabla_c A_{abde}=0.
\]
将$a,d$用逆度量缩并。利用对称性，$g^{ad}A_{cade}=\Ric_{ce}$、$g^{ad}A_{abde}=-\Ric_{be}$，所以
\[
 \nabla^d A_{bcde}+\nabla_b\Ric_{ce}-\nabla_c\Ric_{be}=0.
\]
再缩并$b,e$；$g^{be}A_{bcde}=\Ric_{cd}$。由于$\nabla$与升降/缩并相容，前两项成为两份$\nabla^d\Ric_{cd}$，第三项为$-\partial_c\mathrm{Scal}$，于是得到第一式。$\nabla g=0$使$\nabla^i(\mathrm{Scal}\,g_{ij})=\partial_j\mathrm{Scal}$，从而第二式成立。整个推导是协变张量缩并，不需要把尚未构造的法坐标当先修。
\end{proof}

\begin{corollary}[Schur引理的维数条件]\label{cor:schur}
若连通$n$维黎曼流形上$n\geq3$，且每点所有截面曲率相同，则此共同值在整个流形上为常数。
\end{corollary}
\begin{proof}
共同值函数为$k=\mathrm{Scal}/(n(n-1))$，所以光滑。由$\Ric=(n-1)kg$及缩并Bianchi，
\[
 (n-1)\dd k=\frac{n(n-1)}2\dd k,
 \qquad (n-1)(n-2)\dd k=0.
\]
维数条件给$\dd k=0$，故$k$局部常；连通性使它全局常。二维时系数为零，不能推出$K$恒定，习题给出反例。
\end{proof}

这里的$G$仅作为几何张量定义；广义相对论中的场方程、Lorentz符号及物理解释留第五部分。常曲率空间的整体分类、曲率比较与测地变分也未包含在本章证明范围内。


\Needspace{15\baselineskip}
\section{黎曼几何习题及详细解答}
习题11.1、11.2、11.3、11.4分别改述并扩展梁书上册第二版第3章原题7、12、13、17，均在印刷85/PDF101；11.2所指的成对交换公式在印刷78/PDF94，解释在印刷79/PDF95。其余4题原创。全部解答按本书的曲率与方向槽位重写，原题号、增加部分及双页码另见\texttt{references/EXERCISES.md}。

\Needspace{12\baselineskip}
\subsection*{习题11.1：欧氏球坐标中的联络与散度}
改述并扩展原题7。对$g=\dd r^2+r^2\dd\theta^2+r^2\sin^2\theta\,\dd\phi^2$，$r>0$、$0<\theta<\pi$，求全部非零Christoffel符号；再推出散度和$\Delta_g(r^2)$。区分坐标分量与单位标架分量。
\begin{proof}[解答]
度量及其逆均对角，非零偏导只有$\partial_r g_{\theta\theta}=2r$、$\partial_r g_{\phi\phi}=2r\sin^2\theta$、$\partial_\theta g_{\phi\phi}=2r^2\sin\theta\cos\theta$。代入Christoffel公式得
\begin{align*}
 \Gamma^r_{\theta\theta}&=-r,&\Gamma^r_{\phi\phi}&=-r\sin^2\theta,\\
 \Gamma^\theta_{r\theta}=\Gamma^\theta_{\theta r}&=r^{-1},&
 \Gamma^\theta_{\phi\phi}&=-\sin\theta\cos\theta,\\
 \Gamma^\phi_{r\phi}=\Gamma^\phi_{\phi r}&=r^{-1},&
 \Gamma^\phi_{\theta\phi}=\Gamma^\phi_{\phi\theta}&=\cot\theta.
\end{align*}
其余分量为零（在特殊角度有些列出的函数也可能取零）。这是同一欧氏平坦联络的不同坐标描述，不因系数非零而有曲率。

行列式$d=r^4\sin^2\theta$，本卡$\sqrt d=r^2\sin\theta$，因此对$X=X^r\partial_r+X^\theta\partial_\theta+X^\phi\partial_\phi$，
\[
 \operatorname{div}X=\frac{\partial_r(r^2X^r)}{r^2}
 +\frac{\partial_\theta(\sin\theta X^\theta)}{\sin\theta}
 +\partial_\phi X^\phi.
\]
若改用单位标架$e_r=\partial_r$、$e_\theta=r^{-1}\partial_\theta$、$e_\phi=(r\sin\theta)^{-1}\partial_\phi$，写$X=\widehat X^re_r+\widehat X^\theta e_\theta+\widehat X^\phi e_\phi$，则$X^\theta=\widehat X^\theta/r$、$X^\phi=\widehat X^\phi/(r\sin\theta)$，得到通常含$1/r$和$1/(r\sin\theta)$的物理分量公式。$\operatorname{grad}(r^2)=2r\partial_r$，所以$\Delta(r^2)=r^{-2}\partial_r(2r^3)=6$，与笛卡儿坐标的$\Delta(x^2+y^2+z^2)$一致。公式不在坐标极点直接使用，但函数及其Laplace值在欧氏空间光滑延拓。
\end{proof}

\Needspace{12\baselineskip}
\subsection*{习题11.2：成对交换需要什么假设}
改述并扩展原题12。只用首二槽、末二槽反对称及代数Bianchi，证明四线性张量$A$的成对交换对称性；再说明只有两种反对称性时为何不足。
\begin{proof}[解答]
用简记$A_{abcd}=A(a,b,c,d)$，Bianchi两次展开给
\[
 A_{abcd}=A_{bcda}+A_{cadb}
 =2A_{cdab}+A_{dbac}+A_{adbc}.
\]
而$A_{dbac}+A_{badc}+A_{adbc}=0$，两种反对称给$A_{badc}=A_{abcd}$，所以$2A_{abcd}=2A_{cdab}$。这完成原题所需代数证明；不使用一个不存在的“首末两槽随意交换”规则。

反例取三维空间的对偶基，$\alpha=\varepsilon^1\wedge\varepsilon^2$、$\beta=\varepsilon^1\wedge\varepsilon^3$，令$A(x,y,z,w)=\alpha(x,y)\beta(z,w)$。它对首二槽、末二槽各反对称，但$A_{1213}=1$、$A_{1312}=0$，没有成对交换。其循环和在$(e_1,e_2,e_3,e_1)$上为$-1$，所以缺失的Bianchi条件确实不成立。该例为代数张量反例，不声称它是某个Levi--Civita曲率。
\end{proof}

\Needspace{12\baselineskip}
\subsection*{习题11.3：球面的完整曲率与号差}
改述并扩展原题13。对半径$a>0$的球面卡$g=a^2(\dd\theta^2+\sin^2\theta\,\dd\phi^2)$，求曲率算子、所有四协变曲率分量、Ricci与标量曲率，并核对缩放律。
\begin{proof}[解答]
非零系数为$\Gamma^\theta_{\phi\phi}=-\sin\theta\cos\theta$、$\Gamma^\phi_{\theta\phi}=\Gamma^\phi_{\phi\theta}=\cot\theta$。常量$a^2$在求逆和求导中抵消，与常数缩放不改联络一致。因为$[\partial_\theta,\partial_\phi]=0$，
\begin{align*}
 R(\partial_\theta,\partial_\phi)\partial_\phi
 &=-\partial_\theta(\sin\theta\cos\theta)\partial_\theta
   +\cos^2\theta\,\partial_\theta
 =\sin^2\theta\,\partial_\theta,\\
 R(\partial_\theta,\partial_\phi)\partial_\theta
 &=\bigl(-\csc^2\theta+\cot^2\theta\bigr)\partial_\phi
 =-\partial_\phi.
\end{align*}
因此按第10章的数组约定，独立非零分量可取$R^\theta{}_{\phi\theta\phi}=\sin^2\theta$、$R^\phi{}_{\theta\theta\phi}=-1$，交换最后两个方向槽取负。降指标时本书专用$\mathcal R$的槽序为方向、方向、输入、内积输入，故
\[
 \mathcal R_{\theta\phi\phi\theta}=\mathcal R_{\phi\theta\theta\phi}=a^2\sin^2\theta,
 \quad
 \mathcal R_{\theta\phi\theta\phi}=\mathcal R_{\phi\theta\phi\theta}=-a^2\sin^2\theta,
\]
其余为零。除以Gram行列式$a^4\sin^2\theta$得$K=a^{-2}$，于是
\[
 (\Ric_{ij})=\operatorname{diag}(1,\sin^2\theta)=a^{-2}(g_{ij}),
 \qquad \mathrm{Scal}=2a^{-2}.
\]
$(1,3)$数组与$(0,2)$的Ricci坐标数组均不含$a$，四协变数组乘$a^2$，截面与标量除以$a^2$，四种行为相容。梁书以余矢量交换导数定义曲率，若比较其原题答案，先用第10章的负号及输入重排换到本书约定，不能只凭同一串指标判断正负。
\end{proof}

\Needspace{12\baselineskip}
\subsection*{习题11.4：二维Einstein张量恒零}
改述并扩展原题17：在二维黎曼流形证明$G=0$，说明它不迫使度量平坦。再判断一维时是否有同样的截面解释。
\begin{proof}[解答]
固定一点，正交归一基中所有曲率由$K=\mathcal R_{1221}$决定。取迹得$\Ric_{11}=\Ric_{22}=K$、$\Ric_{12}=0$，所以$\Ric=Kg$、$\mathrm{Scal}=2K$，从而$G=Kg-Kg=0$。本书模型公式也给同一结论。半径$a$的球面$K=a^{-2}>0$，即为$G=0$却不平坦的例子。

原题允许广义黎曼度量；上述题面限定正定，另说明非退化不定情形的代数仍成立：二维空间的交替二形式空间一维，两对槽反对称使$A$为其张量平方的倍数；与$g$的常曲率模型比对得$\Ric=kg$、$\mathrm{Scal}=2k$，故$G=0$。这只补充逐点代数，不在此开展Lorentz几何。一维时方向反对称直接使$R=0$，故$G=0$，但不存在二维截面，不能写“所有截面曲率等于某个唯一的$K$”来解释。
\end{proof}

\Needspace{12\baselineskip}
\subsection*{习题11.5：相容联络的非唯一性}
原创。给定黎曼度量$g$及三形式$H$，以$g(A(X,Y),Z)=H(X,Y,Z)/2$定义$A$，令$\widehat\nabla_XY=\nabla_XY+A(X,Y)$，其中$\nabla$为Levi--Civita联络。证明$\widehat\nabla g=0$、$\widehat T(X,Y)=H(X,Y,\cdot)^\sharp$。在$\R^3$取$H=2\dd x\wedge\dd y\wedge\dd z$，计算一个非零挠率值。
\begin{proof}[解答]
$A$为光滑$(1,2)$张量，故加到联络后仍满足联络规则。两相容性式相减得
\[
 (\widehat\nabla_Xg)(Y,Z)
 =-\tfrac12H(X,Y,Z)-\tfrac12H(X,Z,Y)=0.
\]
又$\nabla$无挠，所以$\widehat T(X,Y)=A(X,Y)-A(Y,X)=H(X,Y,\cdot)^\sharp$。在例子中$A(\partial_x,\partial_y)=\partial_z$、$A(\partial_y,\partial_x)=-\partial_z$，故$\widehat T(\partial_x,\partial_y)=2\partial_z\ne0$。这表明仅“与度量相容”不能保证唯一；基本定理的无挠条件不可删。
\end{proof}

\Needspace{12\baselineskip}
\subsection*{习题11.6：共形变化与二维Schur反例}
原创。设$\widetilde g=e^{2f}g$，证明
\[
 \widetilde\nabla_XY=\nabla_XY+X(f)Y+Y(f)X-g(X,Y)\operatorname{grad}_g f.
\]
在平面$g_0=\dd x^2+\dd y^2$上推出$K_{e^{2f}g_0}=-e^{-2f}(f_{xx}+f_{yy})$；取$f=x^2$说明二维Schur结论失败。
\begin{proof}[解答]
记右端的差为$B(X,Y)$。它对$X,Y$对称，故新联络无挠。又
\begin{align*}
 g(B(X,Y),Z)&=X(f)g(Y,Z)+Y(f)g(X,Z)-Z(f)g(X,Y),\\
 g(Y,B(X,Z))&=X(f)g(Y,Z)+Z(f)g(X,Y)-Y(f)g(X,Z).
\end{align*}
两式相加为$2X(f)g(Y,Z)$，正好补偿$X(e^{2f})$项，所以新联络与$e^{2f}g$相容，由唯一性即为$\widetilde\nabla$。

对平面记$a=f_x,b=f_y$，得到$\Gamma^1_{11}=a$、$\Gamma^2_{11}=-b$、$\Gamma^1_{12}=b$、$\Gamma^2_{12}=a$、$\Gamma^1_{22}=-a$、$\Gamma^2_{22}=b$，下槽交换项相同。计算$R(\partial_x,\partial_y)\partial_y$的第一分量：导数项为$-a_x-b_y$，二次项分别是$-a^2+b^2$及$-b^2+a^2$而抵消；第二分量的导数$b_x-a_y=0$及二次项也抵消。因此
\[
 R(\partial_x,\partial_y)\partial_y=-(f_{xx}+f_{yy})\partial_x,
 \qquad K=-e^{-2f}(f_{xx}+f_{yy}).
\]
取$f=x^2$得到$K=-2e^{-2x^2}$，随点变化。每点当然只有一个二维截面，故“同一点各截面相同”在二维自动成立，却不能推出全局常数。作为号的再检查，$f=-\log y$、$y>0$给$K=-1$，而常数$f$给$K=0$。
\end{proof}

\Needspace{12\baselineskip}
\subsection*{习题11.7：同一联络不决定同一度量}
原创。在连通流形上，证明$g$的Levi--Civita联络也与$h=ug$相容（$u>0$光滑）当且仅当$u$为常数。再说明一般$h$若不限制为共形于$g$，即使共用Levi--Civita联络也未必是常数倍。
\begin{proof}[解答]
$\nabla(ug)=\dd u\otimes g$，因$\nabla g=0$。若它为零，对任意方向$X$及非零$Y$求值给$X(u)g(Y,Y)=0$，所以$\dd u=0$。函数局部常，连通时为一个正常数。反向由常数缩放的结论立即成立；不连通时只能逐分支常数。

在$\R^2$上取$g=\dd x^2+\dd y^2$、$h=4\dd x^2+\dd y^2$。二者的笛卡儿分量都为常数，故同为标准平坦联络；但若$h=cg$，比较两个对角分量要求$c=4$且$c=1$，矛盾。两者单位椭圆正是图\ref{fig:metric-unit-balls}所示。因此“给定度量选出唯一联络”的逆命题不成立。
\end{proof}

\Needspace{12\baselineskip}
\subsection*{习题11.8：标量零、Ricci零与平坦的区别}
原创。在$S^2(1)\times\mathbb H^2$取乘积度量，第二因子用上半平面的$g_H=y^{-2}(\dd x^2+\dd y^2)$。求标量曲率，并判断该例是否Ricci平坦。再在四维欧氏向量空间给出非零代数曲率张量且Ricci迹为零，说明仅从曲率对称性和零Ricci不能推得全曲率为零。
\begin{proof}[解答]
上题的共形公式给第二因子$K=-1$，第一因子$K=1$。乘积Ricci在两个因子上分别为$g_S$及$-g_H$，混合项为零，故标量$2-2=0$，但Ricci并不为零。这先排除把“标量零”误当“Ricci零”的错误。

再取正交基$e_1,\ldots,e_4$，定义对称的平面系数
\[
 k_{12}=k_{34}=1,\quad k_{13}=k_{24}=-1,\quad k_{14}=k_{23}=0.
\]
令$A_{ijji}=k_{ij}$（$i<j$），由两种反对称及成对交换确定同一无序对的其他值；当两个无序指标对不同时令其为零。这给一个四线性张量。对Bianchi，四个指标互异时各项为零；前三个互异而第四重复其中一个时仍各项为零；前三个有重复时两非零项由反对称抵消。因此它满足全部代数曲率对称性。

其Ricci非对角项为零，对角项为$\sum_{j\ne i}k_{ij}$，逐行均为$1-1+0=0$，所以Ricci张量为零；但$A_{1221}=1$，故$A\ne0$。此题只构造逐点代数张量，没有构造Ricci平坦的流形，也不以此代替相应局部实现或整体存在定理。几何实例属于后续专题；本题已经证明不能仅靠这些代数恒等式推出平坦。
\end{proof}

